\section{结论}

本文围绕 Web 搜索与检索增强人工智能的发展脉络，对相关技术和代表性论文进行了梳理。Web 信息检索从爬虫、倒排索引和关键词匹配起步，随后通过 BM25、链接分析和学习排序改善候选结果。预训练语言模型推动了稠密检索与神经重排序，RAG 则把检索结果转化为生成模型的外部证据，使研究目标由返回相关文档扩展到组织证据并生成可核验的回答。这一过程表明，生成框架并未取代检索基础，索引、召回与排序质量仍会直接限制最终答案。

对 WWW 2026 相关论文的分析显示，当前研究正在从单一相关性优化转向可靠性、效率、多模态和智能体决策的联合考虑。系统既要处理检索噪声、证据冲突和生成幻觉，也要避免固定检索深度与过长上下文造成的资源浪费。图像、音频和视频使证据定位扩展到空间、时间和模态对齐，智能体搜索则要求系统继续决定是否搜索、如何改写查询、选择何种工具以及何时停止。因此，检索的质量标准已由“结果相关”扩展为“证据充分、使用可靠且成本可控”。

CFVBench 进一步说明，成功召回相关视频并不意味着模型已经理解其中的细粒度证据，短暂画面、图表数值、界面状态和跨时间片段关系仍可能丢失。AVR 通过判断信息充分性、自适应补帧以及按需调用 OCR 和目标检测改善了视觉召回与回答质量，证明条件计算和专用工具能够弥补部分感知缺口；但残余幻觉、推理成本、自动评价偏差和数据版本差异仍未消除。未来系统应联合优化证据召回、细粒度定位、工具选择与忠实生成，并同时报告效果、延迟、资源消耗和复现条件，在准确性、成本与可验证性之间取得平衡。
